
\section{Théorèmes d'existence, d'unicité et de compatibilité des constantes généalogiques}

La construction détermine un secteur autonome de constantes intrinsèques.
Ses compositions, théorèmes croisés et certificats forment une chaîne
probatoire autonome. L'ordre directeur est

{\small
\begin{equation}
\APP\longrightarrow\TRIT\longrightarrow\TPK
\longrightarrow\text{réalisation fonctionnelle}
\longrightarrow\text{invariant}
\longrightarrow\text{section}
\longrightarrow\text{valeur terminale}.
\label{eq:final-causal-chain}
\end{equation}
}

Avec cet ordre, on obtient cinq clôtures principales.

\begin{enumerate}
\item Le vide électromagnétique cesse d'être une collection de quatre
formules : c'est un opérateur positif à deux feuillets dont le spectre détermine
\(\varepsilon,\mu,Z,c\), avec une amplification nonadique compatible.

\item La criticité cesse de commencer aux décimales de Feigenbaum : la
dodécaphase annule \(\pi\), le mode \(30\) produit \(899\), le germe se clôt
trigonométriquement et entre de manière démontrée dans la classe analytique
unimodale appropriée. L'hyperbolicité universelle est isolée comme un
certificat distinct.

\item La gravité cesse de n'être qu'une carte décimale : la périodicité temporelle de \(108\) phases sélectionne une
section de \(L_G\), le projecteur de rang cinq fixe le support tensoriel et
la réduction TT distingue ce support des deux hélicités sans masse.

\item Le changement hexade--octade acquiert trois réalisations liées :
quantum
aréal, mémoire modulaire et réflexion de mélange. Le hamiltonien holographique
modulaire est central dans cette interface de confluence, mais n'absorbe pas le vide, la criticité,
la gravité, l'arithmétique ni la cosmologie dans un titre artificiellement
unique.

\item Les constantes moins usuelles incluent Catalan, Lévy, Lochs et
Euler--Kronecker. Bendersky--Glaisher, les lacunes et les réalisations cosmologiques
apparaissent également comme des valeurs d'opérateurs et de dynamiques, et non comme une liste
à laquelle on cherche des coïncidences rétrospectives.
\end{enumerate}

\section{Caractère théorématique du corollaire métrologique}

Le \cref{cor:metrological-confluence} et la
\cref{tab:metrological-confluence} constituent la clôture probatoire du
développement. Ce ne sont pas un résumé abrégé du matériau précédent. Pour qu'une ligne
entre dans ce corollaire, elle doit conserver simultanément
\[
\text{antécédent}\mid\text{équation}\mid\text{invariant}
\mid\text{signification}\mid\text{valeur terminale}.
\]
Si l'on élimine l'une quelconque de ces cinq composantes, la ligne cesse d'être une
sortie HMT et devient une simple équivalence numérique. Le statut et
le critère de réfutation ne sont pas des colonnes de substitution : ils auditent de l'extérieur la force et le
domaine du quintuplet complet.

Cette formulation produit trois conséquences structurelles et mathématiques.
Premièrement, la valeur terminale peut être une expression exacte dans une droite de
grandeur ; il n'est pas exigé de la réduire au SI pour la reconnaître. Deuxièmement, un antécédent
peut admettre plusieurs réalisations fonctionnelles ---par exemple, \(T\)
détermine aussi bien la réponse du vide que \(D_A\)--- sans que ces
réalisations soient identifiées.
Troisièmement, un entier partagé comme \(30\), un rapport partagé comme \(4/3\)
ou un scalaire partagé comme \(\log3\) n'établit une confluence que lorsque l'on
exhibe l'application qui le transporte et l'invariant qu'il conserve.

En particulier, l'atlas permet de lire deux équations de clôture qui
restaient auparavant dispersées :
\begin{equation}
\boxed{E_P=k_B\Theta_{\rm clk}\log3},
\qquad
\boxed{
(6\to8)\longrightarrow
\bigl(\gammaH\ell_P^2,\,-\log\rho_A,\,R_{\rm CKM}\bigr).}
\label{eq:status-two-confluences}
\end{equation}
La première fait confluer quantum d'énergie, température et information. La
seconde détermine, à partir d'une interface de confluence commune, un quantum
géométrique, un générateur
modulaire et une réflexion de mélange. Aucune équation n'affirme que ses
composantes sont le même objet ; elle affirme que leurs généalogies ont été
conservées jusqu'au point exact de ramification.

Le contrôle négatif est également immédiat. Si \(E_P\) est utilisé pour sélectionner
\(k_B\), si CKM est utilisé pour recalibrer \(\gammaH\), si \(\log3\) est confondu
avec une matrice de densité ou si l'égalité de deux valeurs décimales remplace
l'entrelaceur, le corollaire ne s'applique pas. L'autonomie de la dérivation consiste en
ce que ces contrôles, comme les dérivations, sont formulés au sein
du développement lui-même.

\section{Stratification logique des identités internes, des réalisations dimensionnelles et des reconnaissances externes}

La force d'une affirmation ne se déduit ni du caractère extraordinaire de sa
conclusion ni de sa proximité avec une valeur connue. Elle se déduit de la chaîne
qui la soutient.

\begin{longtable}{p{36mm}p{96mm}}
\toprule
Statut & Contenu du développement\\
\midrule
\emph{Théorème interne exact}
& Opérateur du vide et ses constitutives ; mode \(30\) ; charge et quotients
quantiques ; échelle thermique ; profil clos du chaos ; réseau
\(3/4/12\) ; identités de Gauss ; sélecteur \(\theta_{108}\) ; famille de
Planck ; réflexion CKM ; loi
\(M_n=(a_n/a_0)^{-6}\) dans sa carte.\\

\emph{Composition mathématique exacte}
& Wien, Stefan--Boltzmann, gaz photonique, Euler--Kronecker,
Bendersky--Glaisher, relation angulaire de Weinberg, reconstruction conjointe
aire--mémoire et autocouplage de Higgs dans l'interface à l'arbre.\\

\emph{Démontré avec structure de départ explicite}
& Résultats qui dépendent d'une carte déclarée, d'une orientation, d'une
mesure, d'un état de bord ou d'une donnée de frontière, sans introduire la valeur
terminale comme sélecteur.\\

\emph{Théorème avec hypothèse explicite}
& Classification de la clôture GNS comme facteur hyperfini
\(III_{\lambda_\partial}\), sous fidélité, compatibilité et persistance
infinie des deux canaux \(90/120\).\\

\emph{Réalisation physique conditionnelle}
& Identification complète de jauge, \(G_F\) renormalisée, réalisation de
Sitter du
rayon CT108, transduction de la mémoire de Perron en densité de spin et
réalisation dynamique du support tensoriel.\\

\emph{Programme ouvert avec critère de réfutation}
& Hyperbolicité complète du renormalisateur de Feigenbaum, valeur propre
sous-dominante de Gauss--Kuzmin--Wirsing et construction d'une action de
Born--Infeld sélectionnée en interne.\\

\emph{Comparaison métrologique externe}
& Cartes SI, comparaisons CODATA/PDG, valeurs classiques des constantes et
classifications opératorielles appliquées uniquement après la construction des poids
HMT.\\
\bottomrule
\end{longtable}

Une interface ouverte ne diminue pas le théorème interne qui la précède. À
l'inverse, un théorème interne ne promeut pas automatiquement une interprétation
physique dont le transducteur n'a pas encore été construit.

\section{Obstructions transversales de typage, de causalité, d'orientation et de compatibilité nonadique}

La construction est soumise à réfutation. Les contrôles suivants
traversent tous les chapitres.

\begin{enumerate}[label=F\arabic*.]
\item \textbf{Typage dimensionnel.} Une identité qui n'est pas homogène dans les
droites \(L_S,L_C,L_T,L_Q\) est écartée, même si elle reproduit une valeur décimale.

\item \textbf{Non-alimentation par la valeur cible.} Aucune valeur SI, CODATA, PDG,
Feigenbaum, GKW, \(H_0\), \(G_F\) ou anomalie magnétique ne peut sélectionner
un opérateur, un coefficient ou une réalisation interne.

\item \textbf{Convention des feuillets.} Inverser silencieusement \(q_+\) et
\(q_-\), ou confondre l'involution du feuillet avec la conjugaison CP, invalide les
parités et les quotients.

\item \textbf{Compatibilité nonadique.} Une fonctionnelle qui change sous
l'amplification \(9^m\), perd l'unité ou ne conserve pas la mémoire du retour
ne possède pas la limite revendiquée.

\item \textbf{Non-fusion par scalaire.} L'égalité de deux valeurs terminales
n'identifie pas leurs domaines. Il faut un entrelaceur qui commute avec les
opérateurs et préserve l'orientation, la retenue et la frontière.

\item \textbf{Séparation fini--infini.} Un certificat jusqu'au niveau huit
ne prouve pas l'hyperbolicité universelle ; un hamiltonien modulaire fini n'est pas
\(-\log\rho_\infty\) dans l'UHF ; une convergence numérique sans borne de
queue ne certifie pas un produit premier.

\item \textbf{Transducteur physique.} La lecture de Sitter, d'Einstein--Cartan,
de Born--Infeld ou de jauge est bloquée si elle ne détermine pas l'objet complet de la
théorie d'arrivée et ses conditions de bord.
\end{enumerate}

\section{Conditions nécessaires pour les extensions mathématiques et les réalisations physiques non équivalentes}

Il ne reste pas de réserve générique du type ``il manque à relier HMT à la
physique''. Il reste des obligations localisées :

\begin{description}[style=nextline,leftmargin=37mm]
\item[CHAOS-OBL-1]
Certifier dans un espace de fonctions analytiques paires le point fixe de
\(\mathcal R\), une unique valeur propre instable, la borne de queue et la
transversalité du germe HMT.

\item[GAUSS-OBL-1]
Fixer l'espace fonctionnel de l'opérateur de Gauss, isoler la valeur propre
sous-dominante par des intervalles et démontrer la stabilité du projecteur sous les
partitions \(9^m\).

\item[HMOD-OBL-1]
Vérifier la compatibilité des états locaux et la persistance avec
densité positive des deux canaux avant d'appliquer la classification
d'Araki--Woods. Ne pas utiliser une densité globale de classe trace inexistante.

\item[EW-OBL-1]
Construire l'entrelaceur entre les feuillets HMT et la base neutre/chargée électrofaible,
et dériver \(\Delta r\) sans employer la valeur expérimentale de \(G_F\).

\item[COS-OBL-1]
Pour de Sitter, démontrer que le réalisateur PCH sélectionne dynamiquement
\(R_{108}\). Pour Einstein--Cartan, construire l'amplitude, le courant, le signe et la
carte dimensionnelle de l'application \(M_n\mapsto s_n^2\).

\item[BORN-OBL-1]
Construire l'action déterminantale de Born--Infeld et démontrer que l'échelle
\(E_*\) est son paramètre, au lieu de l'identifier par les dimensions.

\item[ANOM-OBL-1]
Clore la provenance locale de \(1000\) et \(54+1\) avant de promouvoir
l'évaluation de \(a_\mu\) en dérivation complète.
\end{description}

Chaque obligation possède un objet, un domaine et un critère de réfutation. Aucune n'autorise à
rouvrir les blocs exacts qui la précèdent.

\section{Délimitation formelle du domaine des dérivations généalogiques}

\begin{itemize}
\item Ne redérive pas \(\pi,e,\varphi,\alpha_{\HMT}\) ; les utilise comme entrées déjà
construites lorsqu'une composition ultérieure l'exige.
\item N'entre pas en concurrence avec la Loi générale des masses et ne recalcule pas la masse électronique.
\item Ne redérive pas \(\gammaH\) ; développe sa signification aréale,
informationnelle et modulaire.
\item N'identifie pas le support à cinq composantes à cinq polarisations
d'un graviton sans masse.
\item Ne confond pas une approximation finie de Feigenbaum avec le théorème
universel ni une échelle dimensionnelle avec une action de Born--Infeld.
\item Ne fait pas du hamiltonien holographique modulaire le nom ou
l'explication exclusive de toutes les constantes. C'est une pièce centrale de
l'interface de confluence aire--information--mélange au sein d'un atlas plus vaste.
\end{itemize}

\section{Capacité reconstructive de la généalogie commune sur les constantes mathématiques et physiques}

La singularité métrologique de HMT--MD ne consiste pas à produire beaucoup de
décimales. Elle consiste en ce que la valeur réelle apparaît à la fin d'une séquence qui
conserve la mémoire. La perméabilité et la permittivité sont les deux valeurs propres
quadratiques d'une même réponse ; \(G\) est la section qui clôt un rayon
de Compton avec un rayon gravitationnel ; \(\log3\) est simultanément contenu de
décision et rapport thermique ; Catalan est un moment du quart de torsion ;
\(\gammaH\) oriente le quantum d'aire ; le hamiltonien modulaire conserve quel
canal a produit l'occupation ; la réflexion CKM conserve la coordonnée
responsable du changement de feuillet ;
l'exposant \(-6\) conserve la mémoire stable en déterminant l'échelle
cosmologique.

\begin{proposition}[Résultat principal]
La signification structurelle démontrée dans ce développement s'exprime
par une règle de réalisation : un même antécédent peut admettre
plusieurs fonctionnelles sans perdre sa généalogie, et chacune détermine
simultanément la signification de la constante et la raison de sa valeur
terminale.
\end{proposition}

L'affirmation finale est donc plus précise que ``HMT reproduit des
constantes''. HMT--MD construit des opérateurs, des caractères, des sections et des
transducteurs dont les constantes sont des évaluations. Là où la chaîne
est close, la valeur et sa signification sont démontrées ensemble. Là où subsiste une
interface, le critère de réfutation permet de continuer sans dégrader ce qui est déjà établi ni
présenter comme théorème ce qui est encore un programme.
